% 来源：OEBPS/Text/01_38.xhtml（OPF spine 顺序；仅第一册）
\cleardoublepage
\BookEntry[罗杰·洛温斯坦《拯救华尔街》]{罗杰·洛温斯坦《拯救华尔街》}

\BookCover{images/书单23拯救华尔街.jpg}{[美] 罗杰·洛温斯坦著| 孟立慧译| 广东经济出版社2009年版}

\begin{BookSummaryBox}
{\centering \textbf{《财富》推荐的75本商业人士必读书之一\\ 美国《商业周刊》年度最佳图书}\par}
\end{BookSummaryBox}

\subsubsection*{这本书谈什么？}

这是一群华尔街投机天才的沉浮录。
\par

1994年，华尔街，一家名为长期资本管理公司的对冲基金成立了。这家基金是天才的“梦之队”：掌门人是被誉为能“点石成金”的华尔街“债券套利之父”的约翰·麦利威瑟；合伙人包括以期权定价模型而荣膺1997年诺贝尔经济学奖的罗伯特·默顿和马尔隆·斯科尔斯以及前美国财政部长兼美联储副主席戴维·马林斯等。
\par

在成立之初，这家公司就毫不费力地说服80名创始投资者每人至少拿出1 000万美元；美林证券在当时也购买了其一大笔股份，用来出售给自己的大客户；瑞士联合银行几乎在一夜之间，就把长期资本管理公司定位为其最大的投资对象。然而，在1998年3\textasciitilde{}9月的短短6个月内，该基金却整整亏掉50亿美元，最后濒临破产。美联储不得不出面召集14家大银行组成财团对其救助。
\par

《拯救华尔街》揭示投资市场上的成功以及隐藏在其背后的风险、傲慢与贪婪。作者将长期资本管理公司的崛起以及陨落一一道来。当天才陨落，席卷华尔街的那一场灾难，到底是谁之过？这一个有史以来最生动的资本市场案例，是一本足以让投资者、基金管理者和监管当局深思的金融巨著。
\par

\subsubsection*{芒格为什么谈起这本书？}

1998年10月14日，芒格在加州圣塔莫尼卡市米拉马尔喜来登酒店向基金会财务总监联合会，发表了题为“一流慈善基金的投资实践”的演讲。在这次演讲中，他谈到基金会和大型企业的管理问题，其中“长期资本管理公司”的对冲基金由于对高负债率的投资方法太过自信，结果破产了：
\par

\begin{quotation}
长期资本管理公司是一家成立于1994年的对冲基金，其创办人是一位声誉极佳的华尔街债券交易员和两位诺贝尔经济学奖获得者。他们开发出一些复杂的数学模型，利用债券套利发财。1998年，在所罗门兄弟公司不再从事债券套利和外国金融恐慌的共同影响之下，长期资本管理公司连续两个月出现亏损，由于负债水平极高——当年年度负债1 250亿美元，该公司很快陷入了困境。不到几个月，这家基金就损失了将近20亿美元的基金。美联储被迫为该基金提供了紧急援助，以免整个美国经济的流动性发生连锁反应。这次大失败提醒金融界要注意流动性风险的严重性。《拯救华尔街》就是一本关于这次事件的图书。
\par
\end{quotation}

\subsubsection*{《拯救华尔街》作者罗杰·洛温斯坦如是说}

“一直到20世纪60年代之前，债券交易一直都是一项非常无聊的工作。投资者通常要通过其所在地银行的信托部门买进债券，以追求固定收益。只要债券发行机构不违约，只要投资者对所购债券没有更多的意见，投资者大都会对购买债券感到满意，很少有投资者会积极地去进行债券交易。”
\par

“1971年，美国政府终于放开了对黄金价格的管制，此后，阿拉伯国家发起了石油禁运。在这种情况下，如果债券持有人依然对债券能够带来稳定收益抱有幻想的话，那么，宾夕法尼亚中央铁路公司的倒闭案，则毫不留情地摧毁了这种幻想。此时，大多数债券持有者都已蒙受了极其惨重的损失，他们再也不对那些固定收益感到心满意足了。于是，各国政府都被迫逐渐地取消了固定利润制，并放弃了对货币的管制。一句话，‘固定关系’的世界不复存在了。”
\par

“于是，几乎是在突然之间，黄豆成为了万众瞩目的焦点，钱也成了炙手可热的商品。各地的期货交易所纷纷为各式各样的金融产品设计新的期货合约。一时间，市场上新的金融衍生工具、新期权及新债券交易则大行其道。而此时，专业的证券经纪人也开始醒了过来，并希望能够在这一市场上大显身手。到了20世纪70年代末，像所罗门兄弟公司这样的证券公司，开始以西德尼·霍莫尔之流绝对无力想象的方式，对债券进行分拆和组合交易。他们先将各种各样不同的抵押品集合到一起，然后，再把它们分拆成比较容易交易的小块证券，在市场上进行交易。”
\par
